\section{Radial realization and its geometric limits}
\label{n:sec-geometry-input}

The angular program is fixed before it is placed at large radii. We now construct
a metric with nonnegative Ricci curvature and the coefficient estimates
needed for exact harmonic continuation. The construction is uniform:
one radial factor must accommodate infinitely many independent controls.

Let $h(w)$ denote the fixed-area gauge from the angular construction, so
$d\operatorname{vol}_{h(w)}=d\operatorname{vol}_{g_0}$. Write $w(s)$ for
the fixed periodic reference program, constant near its initial phase.
Extend it on $-2<s<0$ to the round value $w=0$, and set $w(s)=0$ for
$s\leq-2$. The extension stays in the same small parameter neighborhood.
Let $\psi_\nu$ and $\chi_\nu$ be the fixed perturbation directions and
smooth phase bumps at successive crossings. The phase-bump supports are disjoint,
have uniformly bounded phase width, and lie where the selected variation
coefficient has one nonzero sign; each bump is one on a smaller crossing
neighborhood. There are only finitely many smooth types of these data.

Fix $0<a<1$ and $a^2<\kappa_*<1$. For a large start parameter $J$, set
\[
 \eta(t)=t^{-3/4},\qquad
 s(t)=4(t^{1/4}-J^{1/4})\quad(t>0),\qquad s'(t)=\eta(t).
\]
If $s(t_\nu)=s_\nu$ is the center of crossing $\nu$, a sequence
$z=(z_\nu)\in[-1,1]^{\mathbb N}$ prescribes
\begin{equation}\label{n:controlled-angular-input}
 w_z(t)=w(s(t))+
 \sum_\nu z_\nu\eta(t_\nu)^{3/4}\chi_\nu(s(t))\psi_\nu,
 \qquad H_z(t)=h(w_z(t)).
\end{equation}
For $t\leq0$ set $H_z(t)=g_0$. This agrees smoothly with the displayed
formula near zero. At most one summand is nonzero at any time, and the
supports have no finite accumulation. Thus every sequence $z$ defines a
smooth angular path. Its area form is the round one at every time.

\begin{proposition}[Uniform radial realization]
\label{n:prop-radial-realization}
Suppose the fixed reference program, including its negative-phase
interpolation, is close enough to round that its Gauss curvature is
uniformly greater than $\kappa_*$. There are a constant
$C>0$ and, after $C$ is fixed, an integer $J_0$ such that for every integer
$J\geq J_0$ there is a smooth function
$f=f_J:[0,\infty)\to[0,\infty)$, independent of $z$, such that
\[
 g_z=dr^2+f(r)^2H_z(\log r)
\]
extends to a complete smooth metric on $\mathbb R^3$ for every
$z\in[-1,1]^{\mathbb N}$. The metric is Euclidean near zero,
$\operatorname{Ric}_{g_z}\geq0$, with strictly positive Ricci curvature
outside a compact set, and $d_{g_z}(0,(r,x))=r$. Moreover
\[
 a\leq f'(r)\leq1,\qquad f'(r)\to a,\qquad b(t):=e^{-t}f(e^t)\to a.
\]
For every $t\geq0.8J$ the following estimates are uniform over the entire
product box of controls, also after any fixed number of angular
differentiations:
\begin{align}
 &\partial_tH_z=O(\eta),\qquad \partial_t^2H_z=O(\eta^2),
   \qquad \operatorname{tr}(H_z^{-1}\partial_tH_z)=0,
   \label{n:angular-time-input}\\
 &b_{tt}+b_t=-C\eta^2,\qquad
   b_t=O(C\eta^2),\quad b_{tt}=O(C\eta^2),\quad
   b-a=O(Ct^{-1/2}). \label{n:radial-tail-input}
\end{align}
In every fixed smooth angular norm,
$H_z(t)-H(s(t))=O(\eta(t)^{3/4})$ in this range, where the reference
$H(s)=h(w(s))$ includes the negative-phase interpolation.
The controlled links stay in a fixed compact subset of a finite-dimensional
smooth family and have $K_{H_z}>\kappa_*$. All implied constants may depend
on $a$, $\kappa_*$, the fixed angular program, and the fixed $C$; they do not depend on
$J\geq J_0$ or on the control sequence.
\end{proposition}

\begin{proof}
We first quantify the angular speed.  On a fixed-width phase bump,
$|t-t_\nu|=O(\eta(t_\nu)^{-1})=O(t_\nu^{3/4})$; therefore
$\eta(t)/\eta(t_\nu)=1+o(1)$ uniformly as $t_\nu\to\infty$.  Since
$\eta'(t)=-(3/4)t^{-7/4}=O(\eta(t)^2)$ and the fixed phase data have
bounded derivatives, differentiating \eqref{n:controlled-angular-input}
gives, uniformly in $z$,
\begin{equation}
 \partial_t w_z=O(\eta),\qquad
 \partial_t^2 w_z=O(\eta^2).
 \label{n:angular-parameter-derivatives}
\end{equation}
The same bounds hold after any fixed number of angular differentiations.
They also hold through the negative-phase interpolation, where the
controls vanish. Smoothness of $h$ on the fixed compact family also gives
$H_z(t)-H(s(t))=O(\eta(t)^{3/4})$ in every fixed spatial norm.  In particular, if
\[
 A=H_z^{-1}\partial_tH_z,\qquad B=\tfrac12 A,
\]
then $A=O(\eta)$, $\partial_tA=O(\eta^2)$, and
$\operatorname{div}_{H_z}A=O(\eta)$ in the corresponding angular tensor
norms.  Differentiating the fixed area form gives
$\operatorname{tr}(H_z^{-1}\partial_tH_z)=0$, so $A$ and $B$ are traceless.
The perturbation amplitudes $\eta(t_\nu)^{3/4}$ tend uniformly to zero as
$J\to\infty$.  Compactness of the reference phase program, including the negative-phase
interpolation in the chosen near-round neighborhood, therefore
ensures
\begin{equation}
 K_{H_z(t)}>\kappa_*
 \label{n:angular-curvature-lower}
\end{equation}
for all $z$ once $J$ is large.  The supports have no finite accumulation,
so $H_z$ is smooth even when the control sequence is not periodic.

We construct $f$ next.  Choose once a smooth function $0\le\theta\le1$, zero on
$(-\infty,1/2]$ and one on $[0.7,\infty)$, and put
$\theta_J(t)=\theta(t/J)$.  Choose a nonnegative smooth
function $\beta$ supported in $(1,2)$ with $\int\beta(r)\,dr=1$.  For a
constant $C$ to be fixed below, the total slope loss in the tail
\[
 f''_{\mathrm{tail}}(r)=
 \begin{cases}
 -\dfrac{C}{r}\theta_J(\log r)(\log r)^{-3/2},
     &r>e^{J/2},\\
 0, &0\le r\le e^{J/2},
 \end{cases}
\]
is at most $2C(J/2)^{-1/2}$.  Thus, after $C$ has been fixed, $J$ can be
taken large enough that this loss is less than $1-a$.  Set
\[
 m_J=1-a-C\int_{J/2}^{\infty}
                 \theta_J(t)t^{-3/2}\,dt>0,
 \qquad
 f''(r)=-m_J\beta(r)+f''_{\mathrm{tail}}(r),
\]
with $f(0)=0$ and $f'(0)=1$.  The tail is understood to be zero where its
cutoff vanishes.  Then $f=r$ near zero, $f''\le0$, and the total loss of
slope is exactly $1-a$.  Hence $a\le f'\le1$, $f$ is positive on
$(0,\infty)$, and $f'$ tends to $a$.

Set $b(t)=f(e^t)e^{-t}$.  Since $b$ is the average of $f'$ over $[0,e^t]$,
$a\le b\le1$ and $b\to a$.  The tail is exact for $t\ge0.7J$, and
direct differentiation gives
\begin{equation}
 b_t+b=f'(e^t),\qquad
 b_{tt}+b_t=e^tf''(e^t)=-C\eta(t)^2.
 \label{n:radial-tail-identity}
\end{equation}
In this range $f'(e^t)=a+2Ct^{-1/2}$.  We spell out the uniformity in the starting index. Put $u=0.7J$.
Concavity gives $|b_t(u)|=|f'(e^u)-b(u)|\le1-a$, and
\[
 b_t(t)=e^{-(t-u)}b_t(u)
       -C\int_u^t e^{-(t-v)}v^{-3/2}\,dv.
\]
For $u\ge3$, the inequality $\log(t/v)\le(t-v)/u$ bounds the
integral by $2t^{-3/2}$. For $J\ge15$ and $t\ge0.8J$,
\[
 t^{3/2}e^{-(t-u)}\le(0.8J)^{3/2}e^{-0.1J}.
\]
This last expression decreases to zero as $J$ increases. Once $C$ is
fixed, one threshold $J_0$ therefore makes the transient at most
$Ct^{-3/2}$ for every $J\ge J_0$. Consequently
$|b_t|\le3C\eta^2$, $|b_{tt}|\le4C\eta^2$, and, by the first
identity in \eqref{n:radial-tail-identity},
$0\le b-a\le5Ct^{-1/2}$. In particular, uniformly on $t\ge0.8J$,
\begin{equation}
 b_t=O(C\eta^2),\qquad b_{tt}=O(C\eta^2),\qquad
 b-a=O(Ct^{-1/2}).
 \label{n:warping-estimates}
\end{equation}
Here and below $C$ is fixed before $J$ tends to infinity.

The first angular motion occurs only after $s(t)$ reaches $-2$, at
\[
 t_{-2}=(J^{1/4}-1/2)^4
       =J-2J^{3/4}+O(J^{1/2})>0.8J
\]
for sufficiently large $J$.  Thus the exact tail and
\eqref{n:warping-estimates} apply throughout the angular-motion region.
Before that region the metric is round-warped.  Its Ricci eigenvalues are
\[
 \operatorname{Ric}_{rr}=-2f''/f\ge0,
 \qquad
 \operatorname{Ric}^{T}
 =\left(\frac{1-(f')^2}{f^2}-\frac{f''}{f}\right)I\ge0,
\]
including the compact slope transition and the tail cutoff.

It remains to check the complete Ricci tensor while $H_z$ moves.  All
quantities in the following calculation are evaluated at $t=\log r$.
Put $q=1+b_t/b$.  In the outward unit normal $\partial_r$, the shape
operator of the sphere $\{r\}\times S^2$ is
\begin{equation}
 S=\frac1r(qI+B),\qquad \operatorname{tr}S=\frac{2q}{r}.
 \label{n:shape-operator}
\end{equation}
Write $\operatorname{Ric}^{T}$ for the Ricci endomorphism on the unit
angular tangent plane.  The normal, Codazzi, and Gauss identities applied
to \eqref{n:shape-operator} give the exact block formulas
\begin{align}
 r^2\operatorname{Ric}_{rr}
   &=-\frac{2(b_{tt}+b_t)}{b}-\operatorname{tr}(A^2)/4,
       \label{n:ricci-radial-input}\\
 r^2\operatorname{Ric}\bigl(\partial_r,X/f\bigr)
   &=b^{-1}(\operatorname{div}_{H_z}B)(X),
       \label{n:ricci-mixed-input}\\
 r^2\operatorname{Ric}^{T}
   &=\left(\frac{K_{H_z}}{b^2}-2q^2+q-q_t\right)I
       +(1-2q)B-B_t. \label{n:ricci-tangent-input}
\end{align}
In \eqref{n:ricci-mixed-input}, $X$ is an $H_z$-unit vector and
$(\operatorname{div}_{H_z}B)_i=\nabla_jB^j{}_i$.  One may also check
\eqref{n:ricci-radial-input} from
$\operatorname{Ric}_{rr}=-\partial_r\operatorname{tr}S-
\operatorname{tr}(S^2)$.  For \eqref{n:ricci-tangent-input}, use
$\operatorname{Ric}^{T}=f^{-2}K_{H_z}I-
\partial_rS-(\operatorname{tr}S)S$.  These expressions include the mixed
block, which cannot be discarded when testing nonnegative Ricci curvature.

The constants in \eqref{n:angular-parameter-derivatives} give numbers
$M_A,M_M<\infty$, depending on the fixed angular program but neither on
$C$ nor on the controls, such that
\[
 |A|\le M_A\eta,\qquad
 |b^{-1}\operatorname{div}_{H_z}B|\le M_M\eta.
\]
Since $q=1+O(C\eta^2)$ and $q_t=O(C\eta^2)$, the tangential block in
\eqref{n:ricci-tangent-input} satisfies
\[
 r^2\operatorname{Ric}^{T}
   =\left(\frac{K_{H_z}}{b^2}-1\right)I
      +O(\eta+C\eta^2).
\]
The choice $\kappa_*>a^2$ leaves positive slack:
\[
 \tau_0:=\frac{\kappa_*}{a^2}-1>0.
\]
Fix, for instance, $\tau=\tau_0/4$.  Because $b\to a$ uniformly after
late start, the last two displays imply
$r^2\operatorname{Ric}^{T}\ge\tau I$ for all sufficiently large $J$,
once $C$ is fixed.  Meanwhile \eqref{n:radial-tail-identity} and
\eqref{n:ricci-radial-input} give
\[
 r^2\operatorname{Ric}_{rr}
   \ge\left(\frac{2C}{b}-\frac{M_A^2}{4}\right)\eta^2,
 \qquad
 |r^2\operatorname{Ric}_{\mathrm{mixed}}|\le M_M\eta.
\]
Choose $C$ so large that
\begin{equation}
 2C>\frac{M_A^2}{4}+\frac{M_M^2}{\tau}.
 \label{n:ricci-schur-choice}
\end{equation}
This choice uses only the fixed angular program and $a\le b\le1$;
it precedes the choice of $J$.  Choose one $J_0$ after $C$ so that every preceding estimate holds
for every integer $J\ge J_0$, and perform the construction for any such
$J$. Then the Schur complement of the tangential block is bounded below by
\[
 r^2\operatorname{Ric}_{rr}
 -r^2\operatorname{Ric}_{rT}
       (r^2\operatorname{Ric}^{T})^{-1}
       r^2\operatorname{Ric}_{Tr}
 \ge
 \left(2C-\frac{M_A^2}{4}-\frac{M_M^2}{\tau}\right)\eta^2>0.
\]
Thus $\operatorname{Ric}_{g_z}>0$ throughout the angular-motion region,
uniformly for all $z$. Before that region, the round-warped formulas also
give strictly positive Ricci curvature wherever the exact tail has begun,
since $f''<0$ there. Hence $\operatorname{Ric}_{g_z}>0$ for
$r\ge e^{0.7J}$, while it is nonnegative everywhere.

Finally, $f=r$ and $H_z=g_0$ near $r=0$, so
the polar metric $g_z$ extends there as the Euclidean metric.  Every
piecewise smooth path from the origin to $(r,x)$ has length at least $r$
by its radial component, while the radial segment has length $r$.
Therefore $d_{g_z}(0,(r,x))=r$.  Closed metric balls about the origin are
compact in these smooth polar coordinates; hence the metric is complete.
\end{proof}

\subsection{Comparison with the Euclidean metric}
The angular neighborhood can also impose a uniform tensor bound. For
any prescribed $0<\delta<1$, use \eqref{n:eq-gauge-closeness} to require
$(1-\delta)g_0\le h(w)\le(1+\delta)g_0$.
Choose the reference program and its negative-phase interpolation in a
smaller neighborhood with positive margin inside this one. Once their
finite smooth data are fixed, increasing $J$ makes every control
perturbation small enough to stay in the prescribed neighborhood,
uniformly over all sequences $z$. The same tensor bound then holds for
$H_z(t)$ at every time; on the round core it holds automatically.

Since $a\le f(r)/r\le1$ for all $r>0$, comparison with
$g_{\mathrm E}=dr^2+r^2g_0$ gives
\begin{equation}\label{n:eq-global-tensor-comparison}
 a^2(1-\delta)g_{\mathrm E}\le g_z\le(1+\delta)g_{\mathrm E}
 \qquad\text{on }\mathbb R^3.
\end{equation}
The bound extends across the origin, where the metrics agree. Integrating
along curves and taking infima gives the corresponding bounds on their
ambient distances. Thus the distortion of the identity map is controlled
by $a$ and $\delta$, independently of the selected controls.

\subsection{Volume, cone limits, and radial sectional curvature}\label{n:geometric-consequences}
For each allowed $J$ and every control sequence, the metric has
asymptotic volume ratio $a^2$, where the normalization is
$\operatorname{AVR}(g)=\lim_{R\to\infty}\operatorname{vol}_g B(0,R)/(4\pi R^3/3)$.  Indeed, its angular
area form is always the round one, so the distance identity in
Proposition~\ref{n:prop-radial-realization} gives
\[
 \operatorname{vol}_{g_z}B(0,R)
   =4\pi\int_0^R f(r)^2\,dr
   \sim \frac{4\pi a^2}{3}R^3.
\]
The value $a^2$ is independent of the controls.

Let $P$ be the period of the reference angular path.  If
$s(t_n)\pmod P\to s_*$ and $t_n\to\infty$, then for
$r_n=e^{t_n}$ the rescaled metrics $r_n^{-2}g_z$ converge smoothly on
each annulus $0<\epsilon\le r/r_n\le R$ to the metric cone
$d\rho^2+a^2\rho^2H(s_*)$.  Indeed,
$s(t_n+\log\rho)-s(t_n)\to0$ uniformly on such annuli, while
$b(t_n+\log\rho)\to a$ and all controlled perturbations vanish there;
the same chain-rule estimates apply to every fixed derivative.  In the
rescaled metric the ball of radial size $\epsilon$ has diameter at most
$2\epsilon$, so the convergence extends to pointed Gromov--Hausdorff
convergence with the cone tip included. Every phase occurs along a
sequence tending to infinity because $s(t)$ is continuous, increasing,
and unbounded.

The reference path has a phase with simple eigenvalues through the chosen
finite band and another with an isolated double eigenvalue in that band.
Its ordered spectra at these phases differ, so at least one continuous
ordered eigenvalue $\lambda_j(H(s))$ is nonconstant. Its image contains a
nontrivial interval. Phases with distinct values of this eigenvalue give
distinct spectra of the scaled links $a^2H(s)$ and hence nonisometric
links. A pointed cone isometry preserves distance to the tip and restricts
on the unit sphere to an isometry of its intrinsic link metric. The
corresponding pointed tangent cones are therefore nonisometric. In
particular, every control sequence has uncountably many pairwise
nonisometric pointed tangent cones at infinity.

Nonnegative Ricci curvature does not make all radial sectional curvatures
nonnegative here.  Choose a phase $s_*$ where $H_s(s_*)$ is not the zero
tensor; such a phase exists because the angular spectrum changes.  At a
point $x$ where $D=H(s_*)^{-1}H_s(s_*)$ is nonzero, $D$ is self-adjoint
and traceless, so it has a positive eigenvalue.  For times $t_n\to\infty$
with $s(t_n)=s_*+nP$, the radial sectional-curvature endomorphism
$\mathcal K_r=-\partial_rS-S^2$ satisfies
\[
 r_n^2\mathcal K_r
  =\frac{C\eta(t_n)^2}{b(t_n)}I
    +(1-2q)B-B_t-B^2
  =-\frac{\eta(t_n)}2D+o(\eta(t_n))
\]
at $x$.  The control contribution to $B$ is
$O(\eta(t_n)^{7/4})=o(\eta(t_n))$, uniformly in the full control
box. This includes differentiating both the bump and the smooth gauge.  Consequently the radial plane
through a positive eigendirection of $D$ has negative sectional curvature
at all sufficiently late members of this sequence.
